\section{Renormalization in perturbation theory}\label{sec:theory}

According to the standard renormalization in local quantum field theories~\cite{Collins:1984xc}, the renormalization of the matrix elements of an operator does not depend on the external states but is only related to 
the short-distance property of the operator itself. Therefore, one 
can study the renormalization property of the operator in perturbative Green's functions, 
or off-shell quark and gluon matrix elements. For LaMET applications to PDFs, 
we are interested in the non-local operator $O_{\Gamma}(z)=\bar{\psi}(0) \Gamma U(0,z) \psi (z)$ in Eq.~(\ref{eq:operator}). There
are two types of divergences: The linear divergence associated with the Wilson link (its self-energy) and the logarithmic divergences associated with the renormalization of the vertices involving 
the Wilson line and light quark. 
The renormalization is multiplicative, and only the linear divergence has a (linear) $z$ dependence. Therefore the renormalized operator is~\cite{Ji:2017oey,Ishikawa:2017faj,Green:2017xeu, Ji:2020brr}
\bea\label{eq:Reoperator}
O_{\Gamma}(z)_{R}=Z_{O}^{-1} e^{\delta \bar{m} z} O_\Gamma(z),
\eea
where $e^{\delta \bar{m} z}$ contains the linear divergence and 
$Z_{O}$ the logarithmic ones. $\delta \bar{m}$ is not uniquely defined apart
from the linear divergence, introducing a subtraction
scheme dependence which affects the $z$-dependence of the renormalized operator~\cite{Ji:2020brr}.

The linearly-divergent mass renormalization can be calculated in perturbation theory. At one-loop order, the result is independent of lattice action~\cite{Chen:2016fxx}, 
\bea\label{eq:delta_m}
\delta \bar{m} = \frac{2 \pi}{3 a}(\alpha_s+{\cal O}(\alpha_s^2)+...) \ .
\eea
The energy scale of $\alpha_{s}$ can be chosen as $1/a$ to match the lattice results,
\bea\label{eq:alpha_s}
\alpha_{s}(1/a, \Lambda_{\rm QCD}) = \frac{2 \pi}{b_{0} \ln[1/(a \Lambda_{\rm QCD})]} \ ,
\eea
where $b_{0}=11-\frac{2}{3}n_f$ (we will take $n_f$=3) is the QCD $\beta$ function at one-loop order with $n_f$ species of fermions. $\Lambda_{\rm QCD}$ is the non-perturbative QCD scale. 
Including higher-orders in
the $\beta$ function will lead to a more complicated expression. Different choices
of energy scale amount to including high-order corrections.
$\Lambda_{\rm QCD}$ and higher-order $\alpha_s$ terms in $\delta \bar{m}$ also depend on the lattice action. 

The perturbation theory does not converge due to infrared renormalons, see from example~\cite{Ji:1995tm,Beneke:1998ui,Bauer:2011ws,Bali:2013pla}. The existence of the renormalons signals a non-perturbative term in $\delta \bar m$ which is independent of $a$, 
\begin{equation}\label{eq:m0}
      \delta\bar m = m_{-1}(a)/a - m_0 \ , 
\end{equation}
where the minus sign is just a convention.
The uncertainty in summing the perturbation series to get $m_{-1}(a)$
is compensated by the same uncertainty in the non-perturbative $m_0$, leaving the total independent of the renormalon uncertainty. 

Additional uncertainty in $m_0$ comes from the subtraction scheme, or equivalently from the matching to the continuum scheme.  
To reduce the subtraction scheme dependence, we require the renormalized lattice correlation function to be consistent with the $\overline{\rm MS}$ result from continuum perturbation theory at short-distances $z$. To be more concrete, we determine $m_0$ by matching the lattice result to the $\overline{\rm MS}$ one 
within a window $a \ll z < z_S$, where $z_S < 1/\mu$ is the point beyond which perturbation theory ceases to work and $\mu$ is a perturbative scale. The condition $a \ll z$ ensures that the discretization effect on lattice results is small. For this special choice $m_{0c}$, the LaMET expansion does not have a linear term in $1/P^z$~\cite{Ji:2020brr}.

At one-loop order and in dimensional regularization, the renormalization factor of the logarithimic divergence is~\cite{Ji:2020ect,Constantinou:2017sej},
\bea\label{eq:ZO}
Z_{O} = 1+\frac{3 C_{F} \alpha_{s}}{2 \pi} \frac{1}{4-d},
\eea
where $C_{F}=4/3$ is the Casimir operator for the fundamental representation of $SU(3)$ and $d$ is the space-time dimension.
One can resum the logarithimic divergence through solving the renormalization group equation,
\bea\label{eq:ZO_RGE}
\frac{d Z_{O}(\epsilon, \mu)}{d \ln (\mu)}=\gamma Z_{O}(\epsilon, \mu),
\eea
where $\epsilon=(4-d)/2$ and $\gamma=-\frac{3 C_{F}}{2 \pi} \alpha_{s}(\mu, \Lambda_{\rm QCD})$ 
is the leading anomalous dimension, which is independent of the regularization method. 
The form of the leading-order solution of Eq.~(\ref{eq:ZO_RGE}) is independent of regularization scheme and, on the lattice, is
\bea\label{eq:ZO_RS}
Z_{O}(1/a, \mu)=\left(\frac{\ln[1 /(a \Lambda_{\rm QCD})]}{\ln[\mu / \Lambda_{\rm QCD}]}\right)^{\frac{3 C_{F}}{b_0}},
\eea
where the bare ultraviolet (UV) cut-off is $1/a$, and the renormalization scale is $\mu$. 

If one considers the contributions from sub-leading logarithms to the anomalous dimension $\gamma$ when solving the renormalization group equation Eq.~(\ref{eq:ZO_RGE}), we obtain~\cite{Ji:1991pr}
\bea\label{eq:ZO_RS_higher}
Z_{O}(1/a, \mu)=\left(\frac{\ln [1 /(a \Lambda_{\rm QCD})]}{\ln [\mu / \Lambda_{\rm QCD}]}\right)^{\frac{3 C_{F}}{b_0}}\left(1+\frac{d}{\ln[a \Lambda_{\rm QCD}]}\right),\nonumber\\
\eea
where $\Lambda_{\rm QCD}$ and the constant $d$ depend on the specific lattice action.
